
\begin{table}[H]
\centering
\caption{Standardized Effect Sizes for Main Outcomes}
\label{tab:sde}
{\footnotesize
\begin{tabular}{llcccccl}
\toprule
Outcome & Spec. & $\hat{\beta}$ & SE & SD($Y$) & SDE & SE(SDE) & Class. \\
\midrule
\multicolumn{8}{l}{\textit{Panel A: Pooled}} \\
Hire rate & CS-DiD (all) & 0.0091 & 0.0057 & 0.0482 & 0.189 & 0.118 & Large pos. \\
Log earnings & CS-DiD (all) & 0.0027 & 0.0056 & 0.1917 & 0.014 & 0.029 & Small pos. \\
\midrule
\multicolumn{8}{l}{\textit{Panel B: Heterogeneous}} \\
Hire rate & CS-DiD (HS$-$) & 0.0113 & 0.0070 & 0.0517 & 0.220 & 0.136 & Large pos. \\
Hire rate & CS-DiD (BA+) & 0.0069 & 0.0044 & 0.0299 & 0.232 & 0.147 & Large pos. \\
\bottomrule
\end{tabular}
}%
\par\vspace{0.5em}
\parbox{\textwidth}{\scriptsize
\textit{Notes:}
\textbf{Country:} United States.
\textbf{Research question:} Whether state-level reductions in maximum UI benefit duration affect re-employment rates differently for less-educated versus more-educated workers.
\textbf{Policy mechanism:} Seven states shortened maximum UI benefit weeks from 26 to 12--23, reducing the subsidized search window and raising the opportunity cost of unemployment.
\textbf{Outcome definition:} Quarterly hire rate from Census QWI (total hires / beginning-of-quarter employment), aggregated at state-education-quarter level.
\textbf{Treatment:} Binary (state adopted a UI duration cut), with dose variation of 3--10 weeks.
\textbf{Data:} Census QWI sex$\times$education panel, 2007--2020, state-education-quarter, 50 states plus DC.
\textbf{Method:} Callaway--Sant'Anna (2021) staggered DiD, never-treated control, state-clustered SEs.
\textbf{Sample:} Private-sector employment, three education groups (HS or less, some college, BA+).
SDE $= \hat{\beta} / \text{SD}(Y)$ where SD($Y$) is the pre-treatment standard deviation.
Classification refers to magnitude, not statistical significance:
Large ($|$SDE$| > 0.15$), Moderate ($0.05$--$0.15$), Small ($0.005$--$0.05$), Null ($< 0.005$).
}
\end{table}
